
\documentclass[a4paper, 12pt]{article}


\usepackage{mathrsfs}
\usepackage{amsmath,amssymb,amsthm,upref,graphicx,mathrsfs}
\usepackage{enumerate}

\usepackage{geometry}
\geometry{left=3.0cm, right=3.0cm, top=3.0cm, bottom=3.0cm}

\newtheorem{thm}{Theorem}
\newtheorem{lemma}[thm]{Lemma}
\newtheorem{corrollary}[thm]{Corollary}
\newtheorem{proposition}[thm]{Proposition}
\newtheorem{remark}[thm]{Remark}
\newtheorem{question}[thm]{Question}
\newtheorem{definition}[thm]{Definition}
%\newtheorem{definition}[thm]{Definition}
%\newtheorem{proof}[thm]{Proof}
\newcommand{\T}{\mathbb{T}}

\usepackage{color}



%========================================================================%

\title
{{\bf Comments on}\\
	Functional Inequalities on the Quantum Boolean Cube with Applications to Quantum Learning Theory}
\author{by Jiao, Luo, Zhao and Zhou}


\begin{document}


\maketitle


Analysis on the dyadic group $\{-1,1\}^n$ (i.e. Boolean cube in the present paper)   and its Banach-space-valued generalization have important applications in harmonic analysis, probability theory, and functional analysis. This paper is mainly concerned with developing its noncommutative theory. A natural noncommutative counterpart of  dyadic group $\{-1,1\}^n$ is the so-called CAR algebra. The analysis on the CAR algebra goes back at least to \cite{AFS1996}, [CL93] and [LP98] .  Another noncommutative counterpart of the  dyadic group $\{-1,1\}^n$ is the matrix algebra $\mathbb{M}_{2^{n}}:=\mathbb{M}_2^{\otimes n}$ (i.e. quantum Boolean cube). This paper focuses on the analytic theory on the  matrix algebra  (or the quantum Boolean cube) $\mathbb{M}_{2^{n}}$ and its applications.

\medskip

I consider that the main contributions of this paper are as follows

\medskip

(1) It is proved that the quantum Ornstein-Uhlenbeck semigroup is holomorphic on $L_p(\mathbb{M}_{2^{n}})$, $1<p<\infty$. In the classical setting, this result was established by Pisier in his remarkable paper [Pi82a]. More precisely, 



 
 \medskip
 
 My suggestions are listed below
 
 \begin{enumerate}
 	\item 
 \end{enumerate}

\begin{thebibliography}{1}


\bibitem{AFS1996}
S. Ayupov,  S. Ferleger, and F.  Sukochev. Isomorphism between the associative and non-associative $L_p$-spaces of type $II_1$ hyperfinite factors. Mathematica Scandinavica (1996): 271-285.




\end{thebibliography}


\end{document}